\begin{examplebox}
\textbf{例1：已知倾斜角求比例。}
椭圆 $x^2/25+y^2/9=1$，右焦点 $F=(4,0)$，$\theta=60^\circ$。
各例均按正文约定，$A$ 在 $\theta$ 方向，$B$ 在反方向。
由 $a=5,b=3,c=4$，
\[
AF=\frac9{5+4/2}=\frac97,\qquad BF=\frac9{5-4/2}=3,
\qquad \boxed{AF/BF=3/7}.
\]
对应坐标 $A=(65/14,9\sqrt3/14)$、$B=(5/2,-3\sqrt3/2)$，代回椭圆均为 $1$。
\par\medskip\noindent
\textbf{例2：由斜率求比例。}
同一椭圆中 $k=1$，故 $\theta=45^\circ$，
\[
\boxed{\frac{AF}{BF}=\frac{5-2\sqrt2}{5+2\sqrt2}
=\frac{33-20\sqrt2}{17}}.
\]
\par\medskip\noindent
\textbf{例3：由比例反求直线。}
椭圆 $x^2/4+y^2/3=1$，$F=(1,0)$，按上述方向规定 $AF/BF=2$。
因为 $e=1/2$，
\[
\frac{1-e\cos\theta}{1+e\cos\theta}=2
\quad\Longrightarrow\quad \cos\theta=-\frac23.
\]
$\theta\in[0,\pi)$，故 $\sin\theta=\sqrt5/3$，$k=-\sqrt5/2$，直线为
\[
\boxed{y=-\frac{\sqrt5}{2}(x-1)}.
\]
若题目不规定端点命名方向，只给无序长度比，则两条关于 $x$ 轴对称的直线都应讨论。
\end{examplebox}
